\documentclass{situationist-preprint}

\title{Temporal Parts and Present Experience}
\msid{PID-11-b}
\msfamily{PID-11}
\msdiscipline{Personal Identity and the Self}
\mstype{Companion manuscript}
\mskeywords{perdurantism; temporal parts; specious present; present experience; wholly present}

\begin{document}
\maketitle

\begin{abstract}
It is objected to perdurantism that we experience ourselves as wholly present now. We reply that present experience is temporally extended, so its subject is an extended part of the person, who has the experience in virtue of that part; and to feel wholly present is to have all of one's current part present, which perdurantism grants.
\end{abstract}

\section{The objection and the reply}

PID-11-a concludes that persons perdure: every person has a temporal part at every time at which it exists (Corollary~3.4). The objection from present experience says that when I attend to my experience I find myself, the whole of me, here now, not a slice of myself; and that a slice could not be the subject of my experience, since it would then be a different subject from me. It gives phenomenological form to a resistance to temporal parts that has been pressed on metaphysical grounds \citep{Thomson1983}.

\begin{proposition}\label{p:ext}
The subject of a present experience is a temporally extended part of the person, and the person is its subject in virtue of having that part.
\end{proposition}

\begin{proof}
Present experience is not instantaneous. Hearing a phrase of a melody, seeing a movement or feeling a pain swell are experiences of change, and they take time; the present of experience, the specious present, has duration \citep[ch.~15]{James1890}. An instantaneous part cannot have such an experience, since there is no instant at which the whole of it occurs; this is Principle~P4 of PID-11-a. So the subject of a present experience is a part of the person extended over the specious present. The person has the experience as she has a cut when her hand is cut: in virtue of a part, whose states are her states at that time.
\end{proof}

The phenomenology of being wholly present is then explained. What is present to me is all of my current part: nothing of me that exists now is absent. That earlier and later parts are not present is not experienced as an absence. Perdurantism therefore agrees with the phenomenology, and denies only an inference from it, that what is present exhausts what I am.

\begin{thebibliography}{99}\small
\bibitem[James(1890)]{James1890} James, W. (1890). \emph{The Principles of Psychology}. Henry Holt.
\bibitem[Thomson(1983)]{Thomson1983} Thomson, J.~J. (1983). Parthood and identity across time. \emph{The Journal of Philosophy}, 80(4), 201--220.
\end{thebibliography}

\end{document}
